
\examquestion{2023}{4}{Raman transition, atomic Bloch equations and pulse-area sensing}{5--6}
The question describes an effective two-photon
Raman coupling between \(\ket g,\ket e\)
mediated by an auxiliary state \(\ket a\).
The effective Hamiltonian is already written
in terms of a cavity photon and a control field;
the off-resonant auxiliary state has been
eliminated in the model.

\subsection*{(a) Remove the fast free oscillations}
Write
\[
\sigma_+=\ket e\bra g,\quad
\sigma_-=\ket g\bra e,\quad
H_0=\hbar\omega_a a^\dagger a+
\frac{\hbar\omega_{eg}}2\sigma_z.
\]
The interaction in the Schr\"odinger picture
contains
\(\hbar g\mathcal E_c(t)
a^\dagger\sigma_-e^{-i\omega_ct}+\mathrm{h.c.}\).
Under the free interaction picture,
\[
a^\dagger\sigma_-\mapsto
a^\dagger\sigma_-e^{i(\omega_a-\omega_{eg})t},
\]
so the coupling acquires
\[
e^{i(\omega_a-\omega_{eg}-\omega_c)t}
=e^{-i\delta t},
\quad \delta=\omega_{eg}-(\omega_a-\omega_c).
\]
A further rotating-frame phase transformation
transfers this time-dependent phase to the
cavity-number term. With a real control
envelope, an equivalent form is
\[
\boxed{H_I=-\hbar\delta\,a^\dagger a+
\hbar g\mathcal E_c(t)(a^\dagger\sigma_-+
a\sigma_+)}.
\]
The sign of the displayed number term is set
by the selected rotating-frame convention;
the physics depends on the \emph{relative}
two-photon detuning.

\subsection*{(b) Derive the equations for \(z,w\)}
For a strongly driven cavity mode replace
\(a\to\alpha\), \(a^\dagger\to\alpha^*\), treating
\(\alpha\) as the effective coherent amplitude
(including any detuning response).
The atomic coupling is
\(H_{\rm at}=\hbar g\mathcal E_c(t)
(\alpha^*\sigma_-+\alpha\sigma_+)\).
With
\[
z=\langle\sigma_z\rangle,\quad
w=\langle\sigma_+\rangle,\quad
w^*=\langle\sigma_-\rangle,
\]
and commutators
\([\sigma_z,\sigma_\pm]=\pm2\sigma_\pm\),
\([\sigma_-,\sigma_+]=-\sigma_z\),
the Heisenberg equation gives
\[
\dot z=\frac i\hbar\langle[H_{\rm at},\sigma_z]\rangle
=\boxed{2ig\mathcal E_c(t)(\alpha^*w^*-\alpha w)},
\]
\[
\dot w=\frac i\hbar
\langle[H_{\rm at},\sigma_+]\rangle
=\boxed{-ig\mathcal E_c(t)\alpha^* z}.
\]
Its conjugate is
\(\dot w^*=ig\mathcal E_c(t)\alpha z\).
The equations supplied in part (b) assume
that the coupling phase can be taken
time-independent in this rotating description.

\subsection*{(c) Pulse-area solution for complete inversion}
The specified envelope is
\(\mathcal E_c(t)=\mathcal E_0e^{-\gamma|t|}\).
The \emph{pulse area} is
\[
\boxed{\mathcal A=
\int_{-\infty}^\infty\mathcal E_c(t)\diff t
=\frac{2\mathcal E_0}{\gamma}}.
\]
Define the new time coordinate
\(\tau(t)=\int_{-\infty}^t\mathcal E_c(t')\diff t'\).
Then \(\diff/\diff t=\mathcal E_c(t)\diff/\diff\tau\).
The equations reduce to
\[
\frac{\diff z}{\diff\tau}
=2ig(\alpha^*w^*-\alpha w),\qquad
\frac{\diff w}{\diff\tau}=-ig\alpha^*z.
\]
Differentiate \(z\) once more and substitute
the equations for \(w,w^*\):
\[
\frac{\diff^2z}{\diff\tau^2}
=2ig[
\alpha^*(ig\alpha z)-\alpha(-ig\alpha^*z)]
=\boxed{-4g^2|\alpha|^2z}.
\]
The initial atom is in \(\ket g\), so
\(z(-\infty)=-1\), \(w(-\infty)=0\).
The resonant solution is
\[
\boxed{z(t)=-\cos[2g|\alpha|\tau(t)]},\quad
P_e(t)=\frac{1+z(t)}2
=\sin^2[g|\alpha|\tau(t)].
\]
After the complete pulse,
\[
\boxed{P_e(\infty)
=\sin^2\!\left(\frac{2g|\alpha|\mathcal E_0}{\gamma}\right)}.
\]
Complete population inversion requires
\[
\boxed{\frac{4g|\alpha|\mathcal E_0}{\gamma}
=(2m+1)\pi,\quad m=0,1,\ldots,}
\]
equivalently \(2g|\alpha|\mathcal A=(2m+1)\pi\).
This is a \emph{\(\pi\)-pulse} condition.

\textbf{Where does detuning enter?} These Bloch
equations have constant-phase coupling and no
separate atomic detuning term. For their exact
area solution, the effective Raman resonance
must be satisfied, \(\delta=0\), or the relevant
detuning must already have been consistently
absorbed into \(\alpha(\delta)\). If finite
two-photon detuning remains in the atomic
Hamiltonian, the Bloch equations gain extra
precession terms and the simple area result
is not sufficient. For a \emph{constant} drive
of strength \(\Omega_R=2g|\alpha|\mathcal E_c\),
one finds \(\Omega_{\rm eff}=
\sqrt{\Omega_R^2+\delta^2}\), with peak
excitation probability
\(\Omega_R^2/\Omega_{\rm eff}^2<1\) for
\(\delta\ne0\). A pulsed detuned problem must
be solved with its actual time-dependent
Hamiltonian. The question supplies no explicit
cavity response \(\alpha(\delta)\), so a
unique nonzero-detuning inversion condition
cannot be inferred from the given data alone.

\subsection*{(d) Use it as a detuning sensor}
Prepare \(\ket g\), apply the control pulse,
then measure \(P_e\) by state-selective
fluorescence or another population readout.
Scan the frequency difference \(\omega_a-\omega_c\)
and fit \(P_e(\delta)\) to find the two-photon
resonance \(\delta=0\). With a calibrated
drive and detuning-dependent cavity amplitude
\(\alpha(\delta)\), the measured population
provides a frequency-to-population transducer.
For a locally linear sensor, choose a bias
at a nonzero response slope:
\(\delta P_e\simeq P'_e(\delta_0)\,\delta\delta\).
Noise, spontaneous emission, finite pulse
bandwidth and fluctuations set the practical
resolution.
\solcheck{Check the initial condition
\(P_e(-\infty)=0\), final inversion
\(P_e(\infty)=1\) for an odd-\(\pi\) pulse,
and population conservation \(P_e+P_g=1\).
Avoid using a resonance-only pulse-area
solution unchanged at finite detuning.}
